% !PRINTED 443
\weberpage{443}
% running head: § 122.  Irreducibilität.  443

%<<CONT
des Körpers $\Omega$ sein (Bd.\,II, § 149). Es ist dann, wenn $t$ eine Variable bedeutet, $\Phi(t)$ ein Divisor der Klassenfunktion $H_m(t)$, und folglich ist auch die Diskriminante von $\Phi(t)$ ein Teiler der Diskriminante von $H_m(t)$.

Ist $\xi$ eine beliebige ganze Zahl des Körpers $\mathfrak{K}$, so ist die Relativspur in bezug auf $\Omega$ (Bd.\,II, § 175):
\[
\mathfrak{S}\left(\frac{\xi}{t-(k)}\,\Phi(t)\right) = \varphi(t)
\]
eine ganze Funktion von $t$, höchstens vom Grade $\nu-1$, deren Koeffizienten \gesperrt{ganze Zahlen} von $\Omega$ sind. Lassen wir also $t$ in $(k)$ übergehen, so folgt:
\weq{3}{\xi = \frac{\varphi(k)}{\Phi'(k)},}
und daraus:
\weq{4}{b\xi = \beta_0 + \beta_1(k) + \beta_2(k)^2 + \cdots + \beta_{\nu-1}(k)^{\nu-1},}
worin $\beta_0, \beta_1, \ldots, \beta_{\nu-1}$ \gesperrt{ganze Zahlen} in $\Omega$ sind, und $b$ eine ganze rationale Zahl bedeutet, für die man die Diskriminante der Gleichung (2):
\[
b = N[\Phi'(k)]
\]
nehmen kann. Die Körperdiskriminante geht dann in $b$ auf, und \gesperrt{keine der in $b$ nicht aufgehenden Primzahlen ist in $\Omega$ oder in $\mathfrak{K}$ durch das Quadrat eines Primideals teilbar} (Bd.\,II, § 174). Die Zahl $b$ hängt nicht von der besonderen Zahl $\xi$ ab.

Eine ganze Zahl $\xi$ ist nur auf eine Weise in der Form (4) darstellbar, weil sonst $(k)$ einer Gleichung von niedrigerem als dem $\nu$ten Grade genügen würde, was der Voraussetzung widerspricht; \gesperrt{und daraus folgt, daß $\xi$ nur dann durch eine ganze Zahl $\alpha$ des Körpers $\Omega$ teilbar ist, wenn alle Koeffizienten $\beta_0, \beta_1, \ldots, \beta_{\nu-1}$ durch $\alpha$ teilbar sind.} Denn stellt man $\xi/\alpha$ in der Form (4) dar, so müssen die Koeffizienten in dieser Darstellung ebenfalls ganze Zahlen sein.

Ersetzen wir $(k)$ in (3) durch eine andere Wurzel der Gleichung (2), so geht $\xi$ in einen konjugierten Wert über, der ebenfalls eine ganze Zahl ist.

2. Nach dem \gesperrt{Fermatschen} Satze ist für jede beliebige Primzahl $p$, wenn $a$ eine ganze rationale Zahl bedeutet:
\weq{5}{a^p \equiv a \pmod{p}.}
